% !PRINTED 652
\weberpage{652}

In \gesperrt{Abel} a somewhat different expression is found for $k_0$,
namely
\[
k_0 = C + B\sqrt{1+e^2} + \sqrt{h\left(1+e^2+\sqrt{1+e^2}\right)},
\]
which arises from (12) when one sets $a = 1:e$ and then replaces $B$ and $h$
by $Be$ and $he^2$. The expression (12) is thus more general in so far
as it also embraces the special case $a=0$, in which the
biquadratic equation splits into two quadratic equations.

Now, in order to represent the root of a metacyclic equation of degree
5, these expressions for $k_0, k_1, k_2, k_3$ are to be substituted into
the formula (9) of \S.\,186. If we assume $g=2$, the
exponents $r_0, r_1, r_2, r_3$ become congruent in turn with
1, 2, 4, 8, hence equal to 1, 2, 4, 3, and there results, when, as in
\S.\,185,
\[
\begin{aligned}
\tau_0 &= \sqrt[5]{k_0}, & \tau_1 &= \sqrt[5]{k_1} \\
\tau_2 &= \sqrt[5]{k_2}, & \tau_3 &= \sqrt[5]{k_3}
\end{aligned}
\]
is set,
\weq{13}{
\begin{aligned}
\xi = A &+ K_0\, \tau_0^3\tau_1^4\tau_2^2\tau_3 + K_1\, \tau_1^3\tau_2^4\tau_3^2\tau_0 \\
        &+ K_2\, \tau_2^3\tau_3^4\tau_0^2\tau_1 + K_3\, \tau_3^3\tau_0^4\tau_1^2\tau_2.
\end{aligned}
}
The coefficients in this expression, $K_0, K_1, K_2, K_3$, are
rational functions of the $k_0, k_1, k_2, k_3$, which are themselves
cyclically permuted by cyclic permutation of the $k$. According to
\gesperrt{Abel} one is to set
\[
K_0 = A_1 + A_2 k_0 + A_3 k_2 + A_4 k_0 k_2
\]
wherein $A_1, A_2, A_3, A_4$ are rational; but this is not
general enough, because between $k_0 k_2$, $k_0$ and $k_2$ there subsists a
linear relation easily derivable from (12). Most simply one
expresses the $k$ in terms of the three radicals
\[
\begin{aligned}
r &= \sqrt{1+a^2}, & \varrho &= \sqrt{h\left(1+a^2+a\sqrt{1+a^2}\right)} \\
\varrho' &= \sqrt{h\left(1+a^2-a\sqrt{1+a^2}\right)}, & \varrho\varrho' &= hr
\end{aligned}
\]

One can express the radical $\varrho'$ linearly in terms of $r\varrho$ and
$\varrho$, as one sees from the formula
\[
hr^2\varrho' = r\varrho\varrho'^2
\]
when one substitutes on the right for $\varrho'^2$ its expression in terms of $r$
and bears in mind that $r^2$ is rational. Thus one obtains, when
instead of the cyclic permutation of the $k$ one takes the permutation
\[
\begin{pmatrix}
r, & \varrho, & \varrho' \\
-r, & \varrho', & -\varrho
\end{pmatrix}
\]
%CONT>>
